\documentclass[11pt,a4paper]{article}
\usepackage[T1]{fontenc}
\usepackage[utf8]{inputenc}
\usepackage[french]{babel}
\usepackage{lmodern}
\usepackage{amsmath,amssymb,mathtools}
\usepackage{geometry}
\geometry{top=22mm,bottom=22mm,left=23mm,right=23mm,headheight=15pt,headsep=6mm}
\usepackage{microtype,enumitem,booktabs,array,fancyhdr,lastpage,needspace}
\usepackage[hidelinks]{hyperref}
\setlength{\parindent}{0pt}
\setlength{\parskip}{6pt}
\setlist[enumerate]{leftmargin=*,label=\textbf{\alph*)},itemsep=8pt,topsep=4pt}
\pagestyle{fancy}\fancyhf{}
\renewcommand{\headrulewidth}{0.3pt}
\fancyfoot[C]{\small\thepage\,/\,\pageref{LastPage}}
\fancyfoot[L]{\scriptsize Version de travail -- 6 septembre 2026}
\setlength{\emergencystretch}{2em}
\hypersetup{pdftitle={MAT0339 -- Série A, examen 2 -- corrige},pdfauthor={Document de travail pédagogique}}
\fancyhead[L]{\small MAT0339 -- Série A -- Examen 2}
\fancyhead[R]{\small CORRIGÉ ENSEIGNANT}
\begin{document}
{\LARGE\bfseries Série A -- Examen 2}\par
{\large Dénombrement, probabilités et algèbre linéaire}\par
\textbf{Corrigé détaillé et barème}\par
Portée : chapitres 7 à 12. Total : 100 points. Durée proposée : 120 min.\par
\small Document réservé à l’enseignant. Accepter toute démarche équivalente correcte. Une erreur de calcul isolée ne doit pas être pénalisée à chaque étape si la suite est cohérente et reste définie. Les points de domaine, de justification et de validation demeurent distincts.\normalsize\par
\vspace{3mm}\hrule\vspace{4mm}
{\large\bfseries Question 1 -- Compter les issues}\hfill\textbf{15 points}\par
\begin{enumerate}
\item \textbf{5 points.} Les positions sont ordonnées : \(26\cdot25\cdot10^2=65\,000\).
\par{\small\textit{Barème : }Identification des choix/répétitions : 2; expression : 2; calcul : 1.}
\item \textbf{5 points.} L’ordre ne compte pas : \(\binom83=56\).
\par{\small\textit{Barème : }Rôle de l’ordre : 2; expression : 2; calcul : 1.}
\item \textbf{5 points.} Six lettres, avec trois A et deux N : \(6!/(3!2!)=60\).
\par{\small\textit{Barème : }Multiplicités exactes : 2; formule : 2; calcul : 1.}
\end{enumerate}
\vfill {\small\textit{Ancrage dans le manuel : sections 7.1--7.5 et 7.9.}}

\newpage
{\large\bfseries Question 2 -- Tirages sans remise}\hfill\textbf{20 points}\par
\begin{enumerate}
\item \textbf{8 points.} Premier niveau : \(P(R_1)=4/7\), \(P(B_1)=3/7\). Après une rouge, les proportions sont \(3/6\) et \(3/6\); après une bleue, \(4/6\) et \(2/6\). Les chemins \(RR,RB,BR,BB\) ont respectivement les probabilités \[\frac27,\frac27,\frac27,\frac17.\] Leur somme vaut \(1\).
\par{\small\textit{Barème : }Arbre et probabilités conditionnelles : 4; quatre produits : 4.}
\item \textbf{4 points.} Les chemins retenus sont disjoints : \(P=\frac27+\frac17=\frac37\).
\par{\small\textit{Barème : }Chemins pertinents : 2; somme : 2.}
\item \textbf{4 points.} Par définition du conditionnement, \(P(R_1\mid B_2)=\frac{P(R_1\cap B_2)}{P(B_2)}=\frac{2/7}{3/7}=\frac23\). Le dénominateur s’obtient en additionnant les chemins qui réalisent l’événement conditionnant.
\par{\small\textit{Barème : }Dénominateur conditionnel : 2; quotient : 2.}
\item \textbf{4 points.} Par symétrie des positions, \(P(R_2)=4/7\). Or \(P(R_2\mid R_1)=\frac36=\frac12\ne\frac47\). Les événements ne sont donc pas indépendants.
\par{\small\textit{Barème : }Critère : 2; valeurs et conclusion : 2.}
\end{enumerate}
\vfill {\small\textit{Ancrage dans le manuel : sections 8.3--8.4 et 8.7--8.9.}}

\newpage
{\large\bfseries Question 3 -- Essais répétés et espérance}\hfill\textbf{15 points}\par
\begin{enumerate}
\item \textbf{4 points.} \(X\) suit une loi binomiale de paramètres \(n=5\) et \(p=0{,}8\). Chaque essai a deux issues, la probabilité de succès est constante, et les essais sont indépendants.
\par{\small\textit{Barème : }Deux issues, nombre fixe, probabilité constante et indépendance : 1 chacun.}
\item \textbf{4 points.} \(P(X=4)=\binom54(0{,}8)^4(0{,}2)=0{,}4096\). Le coefficient binomial compte les positions des succès.
\par{\small\textit{Barème : }Coefficient : 1; puissances : 2; résultat : 1.}
\item \textbf{3 points.} Le complément est « aucun succès » : \(P(X\ge1)=1-(0{,}2)^5=0{,}99968\).
\par{\small\textit{Barème : }Complément : 1; calcul : 2.}
\item \textbf{4 points.} \(E(G)=6(0{,}8)-4(0{,}2)=4\) dollars par essai. C’est une moyenne théorique à long terme, non un gain garanti à chaque essai.
\par{\small\textit{Barème : }Gains signés : 1; moyenne pondérée : 2; interprétation : 1.}
\end{enumerate}
\vfill {\small\textit{Ancrage dans le manuel : sections 8.5--8.6 et 8.9.}}

\newpage
{\large\bfseries Question 4 -- Angle et projection}\hfill\textbf{15 points}\par
\begin{enumerate}
\item \textbf{3 points.} Par Pythagore, \(\|\mathbf u\|=5,\quad\|\mathbf v\|=4\).
\par{\small\textit{Barème : }Normes et racines : 3.}
\item \textbf{5 points.} \(\mathbf u\cdot\mathbf v=12\). Ainsi \(\cos\theta=3/5\), et \(\theta\approx 53{,}13^\circ\), dans \([0,180^\circ]\).
\par{\small\textit{Barème : }Produit : 2; quotient de cosinus : 2; angle et mode degré : 1.}
\item \textbf{7 points.} \[\operatorname{proj}_{\mathbf v}\mathbf u=\frac{\mathbf u\cdot\mathbf v}{\mathbf v\cdot\mathbf v}\mathbf v=\frac{12}{16}(4,0)=(3,0).\] Le résidu est \((0,4)\). Son produit scalaire avec \(\mathbf v\) vaut zéro; la projection est bien parallèle à \(\mathbf v\).
\par{\small\textit{Barème : }Formule : 2; projection : 2; résidu : 1; orthogonalité : 2.}
\end{enumerate}
\vfill {\small\textit{Ancrage dans le manuel : sections 9.3--9.4 et 9.7--9.8.}}

\newpage
{\large\bfseries Question 5 -- Composer et inverser des matrices}\hfill\textbf{20 points}\par
\begin{enumerate}
\item \textbf{6 points.} Par produit ligne-colonne, \[AB=\begin{pmatrix}2 & 5\\1 & 3\end{pmatrix},\qquad BA=\begin{pmatrix}4 & 3\\1 & 1\end{pmatrix}.\] Les matrices ne commutent pas.
\par{\small\textit{Barème : }Deux produits : 5; comparaison : 1.}
\item \textbf{6 points.} \(\det A=1\ne0\). La formule de l’inverse donne \[A^{-1}=\begin{pmatrix}1 & -1\\-1 & 2\end{pmatrix}.\] Le produit \(AA^{-1}\) est \(I_2\).
\par{\small\textit{Barème : }Déterminant : 2; inverse : 3; produit identité : 1.}
\item \textbf{8 points.} En appliquant l’inverse, \[\mathbf x=A^{-1}\begin{pmatrix}7\\4\end{pmatrix}=\begin{pmatrix}3\\1\end{pmatrix}.\] Vérification : \(2\cdot3+1\cdot1=7\) et \(1\cdot3+1\cdot1=4\).
\par{\small\textit{Barème : }Méthode : 2; solution : 4; deux substitutions : 2.}
\end{enumerate}
\vfill {\small\textit{Ancrage dans le manuel : sections 11.3--11.8 et 12.4--12.6.}}

\newpage
{\large\bfseries Question 6 -- Intersection et classification d’un système}\hfill\textbf{15 points}\par
\begin{enumerate}
\item \textbf{8 points.} La substitution donne \(3+2t=5\), donc \(t=1\) et le point \((3,1,1)\). Le produit de la normale par le directeur vaut \((1,1,1)\cdot(2,1,-1)=2\ne0\), donc l’intersection est unique.
\par{\small\textit{Barème : }Substitution : 3; paramètre et point : 3; unicité : 2.}
\item \textbf{7 points.} La somme des deux premières équations impose \(2x+2z=4\). Si \(k\ne4\), le système est incompatible. Si \(k=4\), la troisième équation est redondante; la différence des deux premières fixe \(y=(3-1)/2\). Avec \(z=s\), toutes les solutions sont \((x,y,z)=(2-s,1,s)\), \(s\in\mathbb R\). Il n’existe aucun cas de solution unique.
\par{\small\textit{Barème : }Condition de compatibilité : 3; classification : 1; famille complète : 3.}
\end{enumerate}
\vfill {\small\textit{Ancrage dans le manuel : sections 10.4--10.7 et 12.3, 12.8--12.9.}}
\end{document}
